\FloatBarrier
\section{Tổng quan báo cáo và phạm vi dữ liệu}

\subsection{Thông điệp chính và mục tiêu phân tích}

Trong phạm vi mẫu được phân công tại mốc 2026-09-16, \textbf{Waterfront Communities-The Island có 161 chỗ ở có giá hợp lệ, gấp 5.0 lần Niagara} ($n=32$), cho thấy đây là khu vực tập trung nguồn cung nổi bật nhất cần được theo dõi trong mẫu; kết quả này không đại diện cho toàn bộ Toronto.

Ở cùng mốc, \textbf{86.75\%} chỗ ở thuộc loại hình nguyên căn và giá trung vị sau kiểm soát chất lượng (QA) đạt \textbf{201.07 CAD/đêm} trên $n=783$ chỗ ở có giá hợp lệ.

Báo cáo hướng tới \textbf{Cơ quan quản lý du lịch thành phố Toronto} nhằm cung cấp cơ sở định lượng cho bốn câu hỏi trọng tâm:
\begin{itemize}
\item Xác định khu vực tập trung nguồn cung cho thuê ngắn hạn và mức độ chênh lệch mặt bằng giá giữa các khu vực.
\item Theo dõi nhịp độ biến động lịch phòng qua 4 mốc chụp.
\item Phân tích điểm mạnh và hạn chế nổi bật trong đánh giá của khách, đồng thời đối chứng mô hình ngôn ngữ lớn (LLM) với mô hình từ điển.
\item Thẩm định các nhận định phân tích để tránh lỗi suy diễn quá mức.
\end{itemize}

\subsection{Giới hạn dữ liệu cần lưu ý}

Dữ liệu được lấy từ \textbf{Inside Airbnb}\footnote{Nguồn dữ liệu và giấy phép: \url{https://insideairbnb.com/get-the-data/}.} (giấy phép CC BY 4.0, không thu thập trực tiếp từ \texttt{airbnb.com}) trên mẫu được giảng viên phân công qua 4 mốc chụp (\text{2025-09-11}, \text{2026-03-19}, \text{2026-06-15}, \text{2026-09-16}), do đó không đại diện cho toàn bộ thị trường thành phố. Riêng mốc cuối (\text{2026-09-16}) không còn quan sát thấy \textbf{72} chỗ ở và \textbf{235} đánh giá có mặt ở các mốc trước (chi tiết tại \text{missing\_ids.csv}). Dữ liệu hiện có không đủ để kết luận nguyên nhân các ID này vắng mặt. Vì vậy, các phép so sánh giữa các kỳ phản ánh mẫu quan sát tại từng mốc, không phải đoàn hệ (cohort) cố định.

\subsection{Quy mô mẫu dữ liệu qua 4 mốc chụp}

Bộ dữ liệu bao gồm ba bảng chính (\text{listings}, \text{calendar}, \text{reviews}) sau bước kiểm soát chất lượng (QA). Quy mô từng bảng qua 4 mốc chụp được tổng hợp trong Bảng~\ref{tab:sample_size}.

\begin{table}[htbp]
  \centering
  \begin{tabular}{lrrr}
    \toprule
    \textbf{Mốc chụp (Snapshot)} & \textbf{Listing sau QA} & \textbf{Calendar sau QA} & \textbf{Review sau QA} \\
    \midrule
    2025-09-11 & 1,000 & 365,000 & 5,000 \\
    2026-03-19 & 1,000 & 365,000 & 5,000 \\
    2026-06-15 & 1,000 & 365,000 & 5,000 \\
    2026-09-16 & 928   & 338,720 & 4,765 \\
    \bottomrule
  \end{tabular}
  \caption{Quy mô mẫu dữ liệu sau QA qua 4 mốc chụp.}
  \label{tab:sample_size}
\end{table}

Ba mốc đầu giữ nguyên đủ \textbf{1,000} chỗ ở và \textbf{5,000} đánh giá; mốc cuối (\text{2026-09-16}) còn \textbf{928} chỗ ở và \textbf{4,765} đánh giá. Các ID đánh giá có thể xuất hiện lại qua nhiều mốc; K5 đã khử trùng theo ID trước khi tổng hợp. Toàn bộ đường dẫn tải và mã băm SHA256 của 12 tệp nguồn được lưu tại \text{download\_manifest.csv}.

% =====================================================================
\FloatBarrier
\section{Giá trung vị giữa 12 khu vực lớn nhất chênh 2.6 lần}

\textbf{Nhận định chính:} Trong 12 khu vực có nhiều chỗ ở có giá hợp lệ nhất tại mốc \text{2026-09-16}, giá trung vị trải rộng từ \textbf{128.61 CAD/đêm} (Dovercourt-Wallace Emerson-Junction) đến \textbf{336.12 CAD/đêm} (Niagara), gấp \textbf{2.6 lần}. Đồng thời, chỗ ở nguyên căn (\textbf{232.88 CAD/đêm}) cũng cao gấp \textbf{2.6 lần} phòng riêng (\textbf{90.5 CAD/đêm}). Do đó, khi đánh giá mức độ đắt đỏ giữa các quận, cần tách riêng yếu tố vị trí khỏi tỷ trọng loại hình nguyên căn, và lưu ý rằng 76/118 khu có $n < 5$ nên chưa đủ cỡ mẫu đại diện.

\begin{itemize}
  \item \textbf{Chênh lệch giá theo khu vực (K1):} Trong 12 khu vực có cỡ mẫu giá hợp lệ lớn nhất tại mốc \text{2026-09-16}, \textbf{Niagara} có giá trung vị cao nhất (\textbf{336.12 CAD/đêm}, $n = 32$), gấp \textbf{2.6 lần} khu vực thấp nhất trong nhóm là \textbf{Dovercourt-Wallace Emerson-Junction} (\textbf{128.61 CAD/đêm}, $n = 24$). Ba khu vực có số lượng chỗ ở có giá hợp lệ lớn nhất là \textbf{Waterfront Communities-The Island} ($n = 161$, trung vị \textbf{279.17 CAD/đêm}), \textbf{Niagara} ($n = 32$, trung vị \textbf{336.12 CAD/đêm}) và \textbf{Annex} ($n = 30$, trung vị \textbf{321.66 CAD/đêm}). Đây là số lượng trong mẫu, không phải mật độ theo diện tích.
  \item \textbf{Chênh lệch theo cơ cấu loại phòng:} Tại mốc cuối, giá trung vị của chỗ ở nguyên căn đạt \textbf{232.88 CAD/đêm}, cao gấp \textbf{2.6 lần} phòng riêng (\textbf{90.5 CAD/đêm}). Do cơ cấu loại phòng khác nhau giữa các quận, chênh lệch giá trung vị giữa các khu vực chịu tác động đồng thời của cả vị trí và tỷ trọng nguyên căn, không nên diễn giải đơn thuần là hiệu ứng nhân quả của riêng vị trí.
  \item \textbf{Giới hạn về cỡ mẫu và đơn vị không gian:} Dự án sử dụng hai cấp độ không gian tách biệt: (1) các KPI từ bảng listing/lịch (K1, K2, K4, K6, K8) nhóm theo khu vực hành chính \text{neighbourhood\_cleansed} --- có \textbf{118} khu có giá hợp lệ sau khi loại giá ngoại lai ở mốc cuối nhưng tới \textbf{76} khu có $n < 5$, chưa đủ cỡ mẫu đại diện nên biểu đồ chỉ xếp hạng 12 khu có cỡ mẫu lớn nhất; (2) bản đồ phân bố và K7 (LLM) nhóm theo 8 vùng phân tích (\text{analysis\_area} A01--A08) để đảm bảo đủ cỡ mẫu văn bản.
\end{itemize}

\reportfigure{fig_space_02_area_price}{Giá đêm trung vị (K1) tại 12 khu vực có cỡ mẫu lớn nhất, mốc 2026-09-16.}

\reportfigure{fig_space_01_listings_map}{Phân bố không gian theo trung vị giá (201.07 CAD/đêm) và 8 vùng phân tích A01--A08.}

% =====================================================================
\FloatBarrier
\section{Tỷ lệ không khả dụng ổn định nhưng cao nhất sát mốc chụp}

\textbf{Nhận định chính:} Tỷ lệ ngày không khả dụng toàn mẫu duy trì ổn định quanh mức \textbf{41.18\%--43.62\%} qua cả 4 mốc chụp, nhưng luôn đạt đỉnh \textbf{57.64\%--71.34\%} ở tháng ngay sát ngày thu thập rồi giảm dần ở các tháng xa hơn trước khi tăng trở lại về cuối cửa sổ lịch 365 ngày. Trong cùng giai đoạn, tỷ lệ yêu cầu thuê dài hạn $\ge 30$ đêm (K4) ổn định quanh \textbf{9.6\%--10.4\%} và tỷ lệ chỗ ở chưa có đánh giá (K6) giảm từ \textbf{5.0\%} xuống \textbf{2.3\%}. Các chỉ số này hữu ích để theo dõi áp lực lịch phòng ngắn hạn nhưng không phản ánh trực tiếp công suất thuê hay doanh thu.

\begin{itemize}
  \item \textbf{Xu hướng khả dụng (K2, K3):} Tỷ lệ ngày không khả dụng toàn mẫu nằm trong khoảng \textbf{41.18\%--43.62\%} qua cả 4 mốc chụp. Khi tách theo từng tháng của lịch 365 ngày tới (K3), tỷ lệ không khả dụng luôn cao nhất ở tháng ngay sát ngày thu thập (\textbf{57.64\%--71.34\%}), giảm dần ở các tháng xa hơn rồi tăng trở lại về cuối cửa sổ. Thời điểm tăng trở lại khác nhau giữa các mốc chụp nên chưa đủ để kết luận về mùa cao điểm.
  \item \textbf{Các chỉ số đi kèm (K4, K6):} Tỷ lệ chỗ ở yêu cầu thuê tối thiểu $\ge 30$ đêm (K4) ổn định quanh 9.6\%--10.4\%, trong khi tỷ lệ chỗ ở chưa có đánh giá (K6) giảm từ 5.0\% xuống 2.3\%.
  \item \textbf{Giới hạn diễn giải bắt buộc:} (1) Ngày \texttt{available = f} gộp chung cả ngày đã được khách đặt lẫn ngày chủ nhà chủ động khóa lịch, do đó \textbf{không được đồng nhất với công suất thuê phòng thực tế hay doanh thu}; (2) Số đánh giá (K5) không phải tổng lượt khách thuê, chỉ tính trên mẫu review được giao, đồng thời tháng biên/quý cuối chưa trọn vẹn số ngày và mốc cuối không còn quan sát thấy một số ID có mặt ở các mốc trước.
\end{itemize}

\reportfigure{fig_time_01_occupancy_by_snapshot}{Tỷ lệ ngày không khả dụng theo tháng lịch (K3) qua 4 mốc chụp.}

\reportfigure{fig_time_02_kpi_trend}{Xu hướng các chỉ số K2, K4 và K6 qua 4 mốc chụp (tách khung chung trục thời gian).}

% =====================================================================
\FloatBarrier
\section{Mô hình từ điển đạt Accuracy cảm xúc cao hơn LLM}

\textbf{Nhận định chính:} Mô hình LLM gán nhãn tích cực áp đảo (\textbf{89.7\%--95.3\%} ở cả 8 vùng phân tích), tập trung khen ở 3 yếu tố: \textbf{chủ nhà}, \textbf{vị trí} và \textbf{tiện nghi}. Đối chứng trên 100 nhãn tay cho thấy \textbf{Baseline từ điển} vượt LLM có ý nghĩa thống kê ở phân loại cảm xúc tổng thể (Accuracy \textbf{0.99 vs 0.91}), trong khi \textbf{LLM} đạt Micro-F1 khía cạnh khen cao hơn trên tập nhãn tay (\textbf{0.8503 so với 0.6934}), nhưng chưa kiểm định ghép cặp cho chênh lệch này. Trong bài làm này, nhóm gọi Gemini API qua \textbf{gói miễn phí (Free Tier, chi phí thực tế 0 USD)}; mức ước tính \textbf{88.09 USD} chỉ nhằm so sánh chi phí tham khảo nếu mở rộng chạy trả phí cho toàn bộ \textbf{729,202} đánh giá của Toronto.

Mẫu LLM gồm \textbf{1,000/1,000} đánh giá hợp lệ: \textbf{955} ở mốc 2026-09-16 và \textbf{45} bổ sung từ mốc 2026-06-15 do các ID này không còn xuất hiện ở mốc cuối. Dữ liệu được trích xuất bằng \texttt{gemini-3.5-flash-lite} (prompt \texttt{v1}, \texttt{temperature=0}, lưu đệm JSONL), kết quả đối chứng với Baseline từ điển trên 100 nhãn gán tay (\texttt{hand\_labels.csv}) được trình bày ở Bảng~\ref{tab:llm_eval}.

\begin{table}[htbp]
  \centering
  \begin{tabular}{ll@{\hspace{16pt}}rrr}
    \toprule
    \textbf{Phương pháp} & \textbf{Thước đo} & \textbf{Giá trị} & \textbf{Khoảng tin cậy 95\%} & \textbf{Cỡ mẫu ($n$)} \\
    \midrule
    LLM      & \texttt{accuracy} & 0.9100 & {[0.8500, 0.9600]} & 100 \\
    LLM      & F1 (macro) & 0.7112 & {[0.3139, 0.8070]} & 100 \\
    BASELINE & \texttt{accuracy} & 0.9900 & {[0.9700, 1.0000]} & 100 \\
    BASELINE & F1 (macro) & 0.9679 & {[0.5799, 1.0000]} & 100 \\
    \bottomrule
  \end{tabular}
  \caption{Đối chứng Accuracy và Macro-F1 giữa LLM và Baseline từ điển trên 100 nhãn tay (bootstrap 95\%).}
  \label{tab:llm_eval}
\end{table}

\reportfigure{fig_llm_02_aspects}{Ba khía cạnh được khen nhiều nhất là chủ nhà, vị trí và tiện nghi trên 1,000 đánh giá hợp lệ.}

\begin{itemize}
  \item \textbf{So sánh định lượng trên 100 nhãn tay (\texttt{hand\_labels.csv}):} Ở phân loại cảm xúc (\texttt{sentiment}), Baseline đạt Accuracy cao hơn LLM có ý nghĩa thống kê (chênh lệch ghép cặp \texttt{llm\_minus\_baseline} = \textbf{-0.0800}, 95\% CI \texttt{[-0.1400, -0.0300]}; LLM lệch \textbf{9} ca, Baseline lệch \textbf{1} ca). Toàn bộ 9 ca lệch của LLM đều rơi vào ranh giới \texttt{mixed\_confusion} giữa \texttt{positive} và \texttt{mixed} khi khách khen tổng thể nhưng kèm một lưu ý nhẹ.
  \item \textbf{Thế mạnh và giới hạn ở bài toán khía cạnh:} Ở trích xuất khía cạnh khen (\texttt{aspects\_praise}), LLM đạt Micro-F1 = \textbf{0.8503}, cao hơn \textbf{0.6934} của mô hình từ điển trên tập nhãn tay; chưa thực hiện kiểm định ghép cặp cho chênh lệch này. Ngược lại, ở khía cạnh bị chê (\texttt{aspects\_complaint}), do tập nhãn tay chỉ có \textbf{7/100} review chứa lời chê và LLM hay gắn nhãn chê cho cả các lưu ý thông tin nhẹ, Micro-F1 của LLM đạt \textbf{0.3636} so với \textbf{0.5333} của Baseline.
  \item \textbf{Giới hạn cỡ mẫu nhãn tay:} Tập đối chứng 100 review lệch lớp mạnh (\textbf{93} \texttt{positive}, \textbf{6} \texttt{mixed}, \textbf{1} \texttt{negative}), khiến khoảng tin cậy 95\% của Macro-F1 rất rộng (\texttt{[0.3139, 0.8070]} với LLM); cần bổ sung thêm mẫu \texttt{negative}/\texttt{mixed} trước khi đưa ra kết luận khái quát hơn.
  \item \textbf{Chi phí thực tế và ước tính tham khảo mở rộng toàn thành phố:} Quá trình xử lý mẫu 1,000 đánh giá trong dự án sử dụng \textbf{Gemini API gói miễn phí (Free Tier)} nên \textbf{không phát sinh chi phí thực tế (0 USD)}. Con số \textbf{88.09 USD} chỉ là \textit{ước tính giả định để so sánh chi phí} nếu chạy trả phí trên toàn bộ \textbf{729,202} đánh giá tại mốc \texttt{2026-09-16} (dựa trên mức đầu vào trung bình 72.88 token/review và đầu ra trung bình \texttt{39.58} token/review theo \texttt{chars/4 from cached results; thinking tokens unknown}, đơn giá tham khảo Standard paid tier \texttt{0.3 USD/1M} token vào và \texttt{2.5 USD/1M} token ra, tra ngày \texttt{2026-11-11} tại \url{https://ai.google.dev/gemini-api/docs/pricing}).
\end{itemize}

% =====================================================================
\FloatBarrier
\section{Thẩm định (Audit) các nhận định phân tích}

\textbf{Nhận định chính:} Trong 5 nhận định điển hình dễ gặp khi đọc nhanh số liệu, không có nhận định nào đạt yêu cầu giữ nguyên (\textbf{4 nhận định phải sửa}, \textbf{1 nhận định chưa đủ căn cứ}). Nguyên nhân chủ yếu không nằm ở phép tính sai mà ở việc suy diễn quá mức từ dữ liệu mô tả sang doanh thu hoặc quan hệ nhân quả.

Nhóm áp dụng quy trình thẩm định 4 bước (Bài 14) --- \textbf{(1) Tính lại số liệu $\rightarrow$ (2) Kiểm tra phương pháp $\rightarrow$ (3) Kiểm tra diễn giải $\rightarrow$ (4) Kết luận: Giữ / Sửa / Chưa đủ căn cứ} --- như tổng hợp trong Bảng~\ref{tab:audit}.

\begin{table}[htbp]
  \centering
  \footnotesize
  \begin{tabular}{p{0.225\linewidth}p{0.065\linewidth}p{0.16\linewidth}p{0.16\linewidth}p{0.29\linewidth}}
    \toprule
    \textbf{Nhận định ban đầu} & \textbf{Số liệu} & \textbf{Phương pháp} & \textbf{Diễn giải} & \textbf{Kết luận \& Bản sửa đúng mức} \\
    \midrule
    1. 86.75\% chỗ ở là nguyên căn, chứng tỏ tiềm năng dòng tiền đầu tư lớn. & Đúng & Đúng & Sai: suy diễn doanh thu/dòng tiền từ số lượng listing. & \textbf{Sửa:} Mốc 2026-09-16 có 928 chỗ ở trong mẫu, trong đó 86.75\% là nguyên căn (không có dữ liệu doanh thu). \\
    \addlinespace
    2. Giá trung bình (289.01 CAD) cao hơn trung vị (201.07 CAD) chứng minh thị trường siêu sang. & Đúng & Lệch: trung bình bị kéo bởi vài ngoại lai cực đoan. & Quá mức: dùng từ ``chứng minh''. & \textbf{Sửa:} Phân phối giá lệch phải; báo cáo dùng trung vị sau lọc IQR (\textbf{201.07 CAD/đêm}, $n = 783$) để phản ánh mức giá điển hình. \\
    \addlinespace
    3. Tỷ lệ không khả dụng 41.18\%--43.62\% cho thấy công suất lấp đầy khách luôn vượt 41\%. & Đúng & Sai: gộp ngày chủ nhà khóa lịch với ngày đã đặt. & Sai: đồng nhất \texttt{available=f} với lượt đặt phòng. & \textbf{Chưa đủ căn cứ} kết luận công suất thuê; chỉ ghi nhận tỷ lệ ngày không khả dụng trên lịch. \\
    \addlinespace
    4. Dovercourt-Wallace Emerson-Junction rẻ hơn Niagara là do kém hấp dẫn về vị trí. & Đúng & Thiếu: chưa kiểm soát cơ cấu loại phòng; 76/118 khu có $n < 5$. & Sai: suy diễn nhân quả đơn biến. & \textbf{Sửa:} Trong 12 khu có cỡ mẫu lớn nhất, trung vị giá dao động từ 128.61 đến 336.12 CAD/đêm và chịu ảnh hưởng của tỷ trọng nguyên căn. \\
    \addlinespace
    5. LLM vượt trội hoàn toàn mô hình từ điển ở mọi tác vụ phân tích đánh giá. & Sai ở cảm xúc & Đúng (bootstrap ghép cặp) & Sai: bỏ qua độ lệch lớp và kết quả thua trên cảm xúc. & \textbf{Sửa:} Baseline thắng có ý nghĩa thống kê trên cảm xúc tổng thể (Accuracy 0.99 vs 0.91); LLM đạt Micro-F1 khía cạnh khen cao hơn trên tập nhãn tay (0.85 so với 0.69), nhưng chưa kiểm định ý nghĩa thống kê. \\
    \bottomrule
  \end{tabular}
  \caption{Kết quả thẩm định 5 nhận định phân tích theo quy trình 4 bước.}
  \label{tab:audit}
\end{table}

% =====================================================================
\FloatBarrier
\section{Phụ lục phương pháp: Kiểm soát chất lượng (QA), định nghĩa KPI và tái lập}

\subsection{Tác động QA lên giá và nguyên tắc xử lý (Bài 10)}

\begin{itemize}
  \item \textbf{Tác động QA lên giá tại mốc cuối (\texttt{2026-09-16}):} Ghi nhận \textbf{114} chỗ ở thiếu giá (giữ nguyên \texttt{NA}, không tự điền). Giá trung vị trước lọc ngoại lai là \textbf{211.09 CAD/đêm} (trung bình \textbf{289.01 CAD/đêm}); sau khi gắn cờ ngoại lai theo quy tắc IQR tính riêng từng \texttt{room\_type} kết hợp trần \texttt{qa.price\_max = 2,000 CAD}, giá trung vị hợp lệ đạt \textbf{201.07 CAD/đêm} ($n = 783$).
  \item \textbf{Nguyên tắc QA theo Bài 10:} Tách giá trị thiếu (\texttt{missing\_price}, \texttt{missing\_rating}, \texttt{missing\_minimum\_nights}, \texttt{missing\_maximum\_nights}) khỏi lỗi ngoài miền (\texttt{invalid\_price}, \texttt{invalid\_rating} ngoài $[1,5]$, \texttt{invalid\_minimum\_nights}, \texttt{invalid\_maximum\_nights}). Giữ dữ liệu gốc và chỉ gắn cờ thay vì tự điền hoặc xóa vật lý. Ngoại lai chỉ là điểm cần kiểm tra, không mặc định là lỗi; từng KPI loại các quan sát không đạt điều kiện của riêng chỉ số đó.
\end{itemize}

\texttt{review\_id\_repeated} ghi nhận \textbf{29,295 lượt ID xuất hiện lại} trong sáu phép so sánh cặp giữa các mốc, không phải 29,295 hàng lỗi. Chỉ báo này có hành động \texttt{record\_only} và được tách khỏi Bảng~\ref{tab:qa_rules}. \texttt{room\_change} ghi nhận 9 lượt thay đổi loại phòng và cũng chỉ mang tính ghi nhận. \texttt{rows\_with\_any\_flag} (4,128) là tổng số hàng có ít nhất một cờ tính cộng dồn theo bảng và mốc, không phải số hàng duy nhất trên toàn bộ tập dữ liệu.

\begin{table}[htbp]
\centering
\begin{tabular}{lr}
\toprule
\textbf{Cờ/chỉ báo QA theo từng bảng và mốc chụp} & \textbf{Tổng lượt/hàng} \\
\midrule
\texttt{rows\_with\_any\_flag} & 4,128 \\
\texttt{comment\_markup} & 3,650 \\
\texttt{missing\_price} & 213 \\
\texttt{price\_outlier} & 169 \\
\texttt{missing\_rating} & 121 \\
\texttt{room\_change} & 9 \\
\texttt{missing\_minimum\_nights} & 1 \\
\texttt{missing\_maximum\_nights} & 1 \\
\bottomrule
\end{tabular}
\caption{Tổng hợp các cờ và chỉ báo QA khác 0 theo từng bảng và mốc chụp; \texttt{review\_id\_repeated} được tách riêng vì là tổng so sánh cặp chồng lặp.}
\label{tab:qa_rules}
\end{table}

\subsection{Định nghĩa 8 chỉ số KPI (K1--K8)}

\begin{table}[htbp]
  \centering
  \footnotesize
  \begin{tabularx}{\linewidth}{lp{0.24\linewidth}p{0.17\linewidth}X}
    \toprule
    \textbf{Mã KPI} & \textbf{Tên chỉ số} & \textbf{Đơn vị} & \textbf{Công thức / Điều kiện chính} \\
    \midrule
    K1 & Giá đêm trung vị & CAD/đêm & \texttt{median(price} hợp lệ\texttt{)} \\
    K2 & Tỷ lệ ngày không khả dụng theo vùng/mốc & \% & \texttt{100*(1-mean(available))} \\
    K3 & Tỷ lệ ngày không khả dụng theo tháng lịch & \% & \texttt{100*(1-mean(available))} theo tháng ngày lịch \\
    K4 & Tỷ lệ yêu cầu thuê dài hạn & \% & \texttt{100*mean(minimum\_nights>=30)} \\
    K5 & Mật độ đánh giá trong mẫu & đánh giá mẫu / chỗ ở / quý & review duy nhất/quý / listing duy nhất trong mẫu \\
    K6 & Listing chưa có review & \% & \texttt{100*mean(number\_of\_reviews==0)} \\
    K7 & Cảm xúc tiêu cực và điểm khía cạnh & \% hoặc \texttt{[-1,1]} & $100N_-/N_{\mathrm{valid}}$; $(P-C)/(P+C)$, với $P$, $C$ là số lượt khen/chê khía cạnh \\
    K8 & Listing khả dụng cao có nhiều review & \% & $100\times\mathrm{mean}(r\geq0.75\land R\geq10)$; $r$ là trung bình \texttt{available} trên ngày lịch hợp lệ; $R$ là \texttt{number\_of\_reviews} \\
    \bottomrule
  \end{tabularx}
  \caption{Định nghĩa chuẩn hóa 8 chỉ số KPI sử dụng trong báo cáo.}
  \label{tab:kpi_defs}
\end{table}

\clearpage
\subsection{Lệnh tái lập và trạng thái bàn giao}

\begin{itemize}
\item \textbf{Chạy toàn bộ pipeline:} \texttt{python run\_pipeline.py --skip-llm}; dùng \texttt{configs/Toronto.yaml} mặc định. Bước LLM đọc cache đã lưu; không cần khóa API. Khi đã có đủ raw hợp lệ, dùng \texttt{--skip-download --skip-llm} để chạy ngoại tuyến. Pipeline sinh bảng và hình, không tự biên dịch PDF báo cáo.
\item \textbf{Kiểm tra tự động:} \texttt{python -m pytest -c notebooks/pytest.ini -q}: \textbf{27 test đạt}. Các công cụ riêng gồm \texttt{check\_report\_numbers.py} và \texttt{check\_security.py}. Checker báo cáo tính lại \textbf{36 giá trị khai báo} từ dữ liệu sạch và dự đoán hiện tại, kiểm chuỗi hiển thị trong cả \texttt{report.md} và \texttt{report.pdf}; không kiểm mọi số, ngữ cảnh nhận xét hoặc bố cục.
\item \textbf{Điều kiện tính toán:} K2 toàn mẫu và K3 giữ ngày lịch hợp lệ kể cả khi thiếu khu vực; phép tổng hợp theo khu vực yêu cầu khu vực xác định. Khi đánh giá LLM, \texttt{parse\_error} không nhận điểm cho cảm xúc, khía cạnh hoặc ngôn ngữ. Tiêu đề hình, ngày review mới nhất và thống kê lỗi được lấy từ dữ liệu lần chạy.
\item \textbf{Tái lập từ raw ngày 22/11/2026:} Chạy đầy đủ pipeline bằng cache trong thư mục tạm, tải mới 12 CSV gzip và 4 GeoJSON; SHA256 của 12 CSV khớp manifest. Pipeline kết thúc không cảnh báo/lỗi; 34 bảng đối chiếu khớp trong sai số $10^{-10}$ (không so timestamp manifest). Mẫu LLM vẫn đủ 1,000 dự đoán hợp lệ; không gọi API mới.
\item \textbf{Làm tròn và provenance:} Annex được hiển thị thống nhất \textbf{321.66 CAD/đêm}, khớp hình sinh lại từ KPI 321.66499999999996. Mã nguồn đang có thay đổi chưa commit; trước khi nộp, ghi đúng \texttt{git rev-parse HEAD} và kiểm \texttt{git status --short}. Không dùng hash cũ của template làm hash bản nộp. Việc đổi mốc 2025-11-11 thành 2026-09-16 là lựa chọn cấu hình của nhóm, cần được giảng viên chấp thuận nếu áp dụng ngoại lệ đề bài.
\end{itemize}
